\documentclass[11pt,a4paper]{article}

% ── Packages (arXiv-safe, no exotic deps) ──────────────────────────────────
\usepackage[utf8]{inputenc}
\usepackage[T1]{fontenc}
\usepackage{lmodern}
\usepackage[margin=1in]{geometry}
\usepackage{microtype}
\usepackage{hyperref}
\usepackage{url}
\usepackage{enumitem}
\usepackage{xcolor}
\usepackage{titlesec}
\usepackage{abstract}

\hypersetup{
  colorlinks=true,
  linkcolor=black,
  citecolor=black,
  urlcolor=blue!50!black,
  pdftitle={OCPP-AI: Substrate-Agnostic Article 50 Transparency Receipts},
  pdfauthor={Veronica S. Dawkins}
}

\titleformat{\section}{\large\bfseries}{\thesection.}{0.6em}{}
\titleformat{\subsection}{\normalsize\bfseries}{\thesubsection}{0.6em}{}

\setlength{\parskip}{0.5em}
\setlength{\parindent}{0pt}

% ── Title block ────────────────────────────────────────────────────────────
\title{
  \vspace{-2em}
  An Open Cryptographic Provenance Protocol for AI System Outputs (OCPP-AI):\\
  Substrate-Agnostic Article 50 Transparency Receipts Compatible\\ with eIDAS 2.0 and the European Blockchain Services Infrastructure
}
\author{
  Veronica S.\ Dawkins\thanks{LedgerProof Foundation (in formation), Delaware 501(c)(3) under Adler \& Colvin counsel and Dutch Stichting EU subsidiary under NautaDutilh counsel. Correspondence: \texttt{veronica@ledgerproofhq.io}. Specification artifacts: \url{https://spec.ledgerproofhq.io}.}
}
\date{3 June 2026}

\begin{document}
\maketitle

\begin{abstract}
\noindent
Article 50 of Regulation (EU) 2024/1689 (the ``EU AI Act'') imposes transparency obligations on providers and deployers of AI systems with respect to natural persons, supervisory authorities, and the public. Operational implementation requires machine-readable evidence that survives the continued cooperation of any single deployer, vendor, or jurisdiction. Existing vendor-proprietary logging mechanisms fail this requirement by structurally coupling enforcement evidence to individual cloud platform operators outside Union jurisdiction.

This paper specifies the Open Cryptographic Provenance Protocol for AI System Outputs (OCPP-AI), a substrate-agnostic protocol for machine-readable cryptographic receipts addressing Article 50 sub-obligations 50(1), 50(2), 50(4), and 50(6). The protocol separates an \emph{Anchor Interface} specification (defining the properties any conforming anchor substrate must satisfy) from substrate-specific reference implementations: Bitcoin \texttt{OP\_RETURN} as the current reference implementation and the European Blockchain Services Infrastructure (EBSI) as the primary enterprise-government alternative.

We demonstrate architectural compatibility with W3C Verifiable Credentials 2.0 through a published JSON-LD context that wraps OCPP-AI receipts as Verifiable Credentials without loss of information, and architectural complementarity with Regulation (EU) 2024/1183 (eIDAS 2.0) through a composition pattern that enables qualified evidence presentation through established Qualified Trust Service Providers without modification of the core receipt format.

The specification preserves GDPR Article 17 erasure mechanics through structural exclusion of personal data at the receipt schema validation layer and at the anchor payload format. Reference implementations exist in Python, TypeScript, and Rust under Apache License 2.0, with twenty-nine framework adapters covering the principal AI orchestration ecosystems. We propose OCPP-AI for stakeholder consultation at the European Commission AI Office, the European AI Board, CEN-CENELEC Joint Technical Committee 21, and the IETF SCITT Working Group.

\medskip
\noindent\textbf{Keywords:} EU AI Act, Article 50, cryptographic provenance, verifiable credentials, eIDAS 2.0, European Blockchain Services Infrastructure, supply-chain transparency, SCITT, GDPR-compatible cryptography.
\end{abstract}

\section{Introduction}

\subsection{The structural problem of Article 50 evidence}

The European Union's Artificial Intelligence Act, Regulation (EU) 2024/1689, becomes applicable on 2 August 2026 for the transparency obligations of Article 50. These obligations require providers and deployers of AI systems to (1) disclose to natural persons that they are interacting with an AI system, (2) mark synthetic media in a machine-readable manner, (3) notify subjects of emotion-recognition and biometric-categorization systems, and (4) disclose AI-generated text on matters of public interest, with a specified exemption for content subject to human editorial review.

Operational implementation of these obligations creates a structural requirement that the existing AI deployment ecosystem does not satisfy: Article 50 evidence must survive the continued cooperation of the deployer, the vendor, and any single jurisdiction's regulatory authority. Evidence may need to be presented to a national competent authority of a Member State other than that of the deployer's seat, to a court considering allegations against a deployer that has ceased to operate, or to a natural person exercising a right of recourse against a provider-deployer combination of which neither remains commercially incentivized to preserve the evidence.

The contemporary practice of compliance logging through vendor-proprietary mechanisms fails this requirement. Such logs are operationally accessible only with the cooperation of the vendor operating them; they are subject to data-retention policies set by the vendor; and they are accessible to supervisory authorities only through processes that depend on the vendor's continuing presence and cooperation in the jurisdiction.

This structural property of vendor-operated logging has a regulatory consequence: the enforcement infrastructure of the AI Act is, in the default deployment pattern, effectively captured by the cloud platform operators on whose infrastructure the AI systems run. The Union's regulatory authority over Article 50 is, in this default, conditional on the continued cooperation of operators outside the Union's jurisdictional reach.

\subsection{The architectural requirement}

The architectural property that escapes this capture is \emph{independent verifiability without the cooperation of the deployer, the vendor, or any single jurisdiction's authority}. A natural person, a supervisory authority, or a court must be able to verify Article 50 evidence using only publicly available substrate access and a published verifier specification. No verifier may be required to obtain credentials from, enter into agreement with, or otherwise depend on the continuing cooperation of any of the parties to the original AI system transaction.

This requirement is satisfied by the architectural pattern of cryptographic attestations published to public, immutable substrates: signed receipts batched into Merkle trees whose roots are published to a substrate that satisfies a defined Anchor Interface specification (immutability under adversarial conditions, public verifiability without authentication, jurisdictional neutrality, demonstrated operational durability). This paper specifies such a protocol.

\section{Related work}

\textbf{Certificate Transparency.} RFC 6962 and RFC 9162 established the technical pattern of Merkle-tree-aggregated publication of signed attestations to publicly verifiable logs. CT logs are a viable substrate for non-EU evidence applications but do not satisfy the jurisdictional neutrality requirement of EU AI Act evidence in the form deployed for X.509 certificate transparency, where CT logs are operated by a small number of commercial certificate authorities. The architectural pattern is preserved; the substrate is changed.

\textbf{IETF SCITT.} The IETF Supply Chain Integrity, Transparency, and Trust (SCITT) Working Group has produced specifications for cryptographic supply-chain attestations, including the SCITT Architecture and SCITT Receipts. OCPP-AI is formally specified as a SCITT profile in IETF Internet-Draft \texttt{draft-dawkins-scitt-ai-article50-00}, currently under review by the SCITT WG.

\textbf{W3C Verifiable Credentials 2.0.} The W3C VC 2.0 Data Model (W3C Recommendation, May 2025) provides a standardized format for machine-readable cryptographic attestations. Architectural compatibility with VC 2.0 is desirable for participation in the broader ecosystem of verifiable credentials including those produced under eIDAS 2.0.

\textbf{eIDAS 2.0 and the European Digital Identity framework.} Regulation (EU) 2024/1183 establishes the European Digital Identity Wallet framework, depending architecturally on W3C Verifiable Credentials and on Qualified Trust Service Providers (QTSPs). OCPP-AI is not itself an eIDAS-qualified protocol; OCPP-AI signatures are not Qualified Electronic Signatures per eIDAS Article 25. The relationship is one of architectural complementarity: OCPP-AI provides the Article 50 evidence layer; eIDAS 2.0 provides the broader qualified-evidence framework.

\textbf{The European Blockchain Services Infrastructure (EBSI).} EBSI, operated by the European Commission and the European Blockchain Partnership Member States, is a permissioned blockchain network designed for cross-border publication of trusted credentials by Member State governments and EU institutions. EBSI is identified by this paper as the primary enterprise-government alternative anchor substrate satisfying the OCPP-AI Anchor Interface for use cases requiring Union-sovereign verification.

\section{The OCPP-AI specification}

\subsection{Receipt structure}

OCPP-AI receipts comprise three structural layers: an envelope establishing chain identity, a content body conforming to one of four Article 50 sub-obligation profiles, and an anchor reference resolving to a substrate-level anchor record. Receipts are encoded using the Concise Binary Object Representation (CBOR; RFC 8949) deterministic encoding subset and signed using Ed25519 (RFC 8032).

The content body profiles are: \textbf{ChatbotSession} (Article 50(1)), \textbf{GeneratedContent} (Article 50(2)), \textbf{HumanReview} (Article 50(4)), and the manner-of-disclosure field present across profiles (Article 50(6)). Profile schemas constrain field values to bounded vocabularies and reject personally identifiable data at the schema validation layer, consistent with GDPR Article 5(1)(c) data minimization.

\subsection{Receipt aggregation and anchoring}

Receipts are aggregated into batches whose Merkle roots are constructed per RFC 6962. Each Merkle root is published to a substrate satisfying the Anchor Interface (Section~\ref{sec:anchor}) using the 36-byte OCPP-AI anchor payload format: a fixed four-byte ASCII prefix \texttt{"LPR1"} followed by the 32-byte SHA-256 Merkle root.

The fixed-length anchor payload format serves three purposes: machine-distinguishability from non-OCPP-AI uses of the substrate, structural exclusion of personal data at the anchor layer, and preservation of substrate-level cost amortization across the receipts in a batch.

\subsection{Verification}

Verification of an OCPP-AI receipt proceeds in six steps: (1) parse the receipt as canonical CBOR, (2) verify the Ed25519 signature against the publisher's published public key, (3) verify schema conformance of the content body to one of the Article~4 profiles, (4) resolve the anchor reference to a substrate-level record, (5) verify the OCPP-AI anchor payload format including the \texttt{"LPR1"} magic prefix, and (6) reconstruct the Merkle proof from the receipt to the anchor's published Merkle root per RFC 6962. No verification step requires authentication to any party.

\section{The Anchor Interface}
\label{sec:anchor}

The Anchor Interface is the abstract specification of the properties any acceptable anchor substrate must satisfy. The interface specifies eight properties, summarized here and specified in full in the companion document at \url{https://spec.ledgerproofhq.io/anchor-interface-v1.html}.

\begin{description}[leftmargin=1.4em]
\item[I-1 Immutability under adversarial conditions.] The substrate must provide a cryptographic guarantee that no party can alter, remove, or substitute the published anchor record without producing a publicly observable inconsistency.
\item[I-2 Public verifiability without authentication.] Any natural or legal person with internet access must be able to verify the presence and content of any published anchor record without identifying themselves to the substrate operator.
\item[I-3 Jurisdictional neutrality.] The substrate's continued operation must not depend on the consent or non-interference of any single sovereign state or single private legal entity; substrate operation must be distributed across at least three independent jurisdictions.
\item[I-4 Demonstrated operational durability.] The substrate must have demonstrated continuous operational availability for at least 36 months prior to anchor publication.
\item[I-5 Deterministic resolution of anchor records.] Any independent verifier must be able to deterministically resolve the anchor record content and substrate position from the anchor reference.
\item[I-6 Bounded anchor payload format.] The substrate must accept the fixed 36-byte OCPP-AI anchor payload format.
\item[I-7 OCPP-AI Anchor Payload Format.] The fixed 36-byte structure \texttt{"LPR1"} $\|$ \texttt{merkle\_root\_32}.
\item[I-8 GDPR Article 17 compatibility.] The anchor payload format must not permit publication of personal data; receipt-level erasure must not invalidate the anchor.
\end{description}

\subsection{Reference implementation: Bitcoin OP\_RETURN}

The Bitcoin mainnet, with anchor payload publication via the \texttt{OP\_RETURN} script opcode, satisfies all eight properties of the Anchor Interface. Bitcoin OP\_RETURN anchoring is appropriate for use cases prioritizing maximum operational durability, jurisdictional neutrality, and independence from any single state's regulatory authority.

\subsection{Alternative implementation: European Blockchain Services Infrastructure (EBSI)}

EBSI satisfies properties I-1 through I-3 and I-5 through I-8, with property I-4 under continuing evaluation pending EBSI's accumulation of the 36-month operational record threshold. EBSI is appropriate for use cases requiring Union-sovereign verification, where the deployer is a public-sector entity, or where Member State competent authorities have indicated a preference for EU-operated anchor substrates.

\subsection{Dual-anchor deployment}

Dual-anchor publication---the same receipt-batch Merkle root published to both Bitcoin and EBSI in parallel---is the recommended configuration for deployers operating under EU jurisdictional reach whose AI system outputs are syndicated across both EU and non-EU jurisdictions. Verifiers may accept evidence from either anchor; verifiers operated by EU institutions may require evidence from the EBSI anchor as a matter of supervisory policy without disadvantaging non-EU verifiers, who may accept the Bitcoin anchor.

\section{Compatibility with W3C VC 2.0 and eIDAS 2.0}

\subsection{W3C VC 2.0 wrapping}

A JSON-LD context published at \url{https://spec.ledgerproofhq.io/contexts/lpr-v1.jsonld} maps OCPP-AI content body fields to W3C VC 2.0 \texttt{credentialSubject} properties. The mapping is one-to-one and reversible; an OCPP-AI receipt may be losslessly serialized as a W3C VC 2.0 Verifiable Credential and recovered from that form. When wrapped as a VC, the receipt's \texttt{issuer} property identifies the publishing deployer; the LedgerProof Foundation does not assert issuer authority over receipts published by deployers.

\subsection{eIDAS 2.0 architectural complementarity}

OCPP-AI signatures are Ed25519 and do not constitute Qualified Electronic Signatures per eIDAS Article 25. The LedgerProof Foundation is not a Qualified Trust Service Provider. The relationship between OCPP-AI and eIDAS 2.0 is therefore one of architectural complementarity rather than overlap: an OCPP-AI receipt, wrapped as a W3C VC 2.0 credential, may be presented through an eIDAS-compliant Qualified Trust Service Provider for use cases requiring qualified evidence. The composition does not modify the OCPP-AI receipt format; the qualification attaches at the wrapping layer.

This complementarity preserves the Union's existing investment in eIDAS infrastructure. OCPP-AI does not duplicate, compete with, or replace eIDAS trust services. It provides the Article 50 cryptographic provenance evidence layer that complements them.

\section{GDPR compatibility}

The OCPP-AI receipt schema is constructed to preclude the inclusion of personal data at the receipt layer. Fields that could otherwise carry personal data (deployer\_id, reviewer\_role, review\_rationale) are bounded to specific identifier formats (Legal Entity Identifier, VAT identifier, bounded role vocabulary). Schema validation rejects email addresses, telephone numbers, government identifiers, and free-text PII patterns.

At the anchor layer, the 32-byte SHA-256 Merkle root carries cryptographically no recoverable information about underlying receipt content. The anchor is GDPR-compatible by structure; no party may publish personal data to the anchor through this protocol.

Receipt-level erasure under GDPR Article 17 is performed at the receipt store layer. Erasure invalidates the Merkle proof for the erased receipt without invalidating proofs for other receipts in the same batch. The anchor is unaffected by erasure. This separation is essential for joint AI Act / GDPR compliance and is a normative architectural commitment of the protocol.

\section{Implementation status and ecosystem}

Reference implementations of OCPP-AI are published under Apache License 2.0 at \url{https://github.com/ledgerproof} in Python, TypeScript, and Rust. Twenty-nine framework adapters cover the principal AI orchestration ecosystems (LangChain, LlamaIndex, OpenAI SDK, Anthropic SDK, Google AI, Vertex AI, AWS Bedrock, Azure OpenAI, Mistral, Cohere, Together, Groq, Hugging Face, AI21, Replicate, xAI, DeepSeek, Qwen, Reka, Voyage, Cerebras, IBM watsonx, Snowflake Cortex, Perplexity, Fireworks, Aleph Alpha, Haystack, Semantic Kernel, and Mistral Codestral).

A command-line verifier and a browser-based verifier (operating with no SaaS dependency) are published under the same license. Conformance test vectors are published at \url{https://spec.ledgerproofhq.io/test-vectors-v1/}.

\section{Governance and standardization}

The LedgerProof Foundation, the entity maintaining the reference implementation, is in formation as a Delaware 501(c)(3) public charity under Adler \& Colvin counsel, with a Dutch Stichting EU subsidiary under NautaDutilh counsel. The specification is published under CC BY 4.0. The Foundation does not assert patent encumbrance over the specification and undertakes that the specification will not be made subject to commercial licensing.

The Foundation contributes to standardization activity at CEN-CENELEC Joint Technical Committee 21 (through national mirror committees including DIN Germany and AFNOR France), at the IETF SCITT Working Group, and at W3C. The Foundation will support the adoption of this specification or its successors by other standards bodies including ISO, CEN-CENELEC, IETF, and W3C if requested.

\section{Discussion: the Brussels Effect and strategic autonomy}

The architectural choices specified in this paper reflect a substantive position on the proper relationship between AI regulation and the infrastructure on which AI systems run. We argue that:

\begin{enumerate}[leftmargin=1.6em]
\item The Union's regulatory authority over Article 50---and therefore the Brussels Effect for AI regulation---depends on the existence of an enforcement infrastructure that the Union can operate without depending on the continuing cooperation of operators outside the Union's jurisdictional reach.

\item The contemporary default of compliance logging through vendor-proprietary mechanisms forfeits this infrastructure. Article 50 enforcement under that default is conditional on hyperscaler cooperation.

\item The architectural pattern that escapes this conditional dependency is cryptographic provenance evidence published to public, immutable substrates satisfying a defined Anchor Interface. The Anchor Interface is substrate-neutral by design; it admits both substrates operated outside Union control (Bitcoin) and substrates operated under Union control (EBSI).

\item The Union's adoption of an architectural pattern that admits both forms of substrate, with EBSI as a first-class implementation, preserves Union strategic autonomy over Article 50 enforcement without sacrificing the operational durability and jurisdictional neutrality that long-term evidence preservation requires.
\end{enumerate}

We do not take a position on which substrate should be preferred in specific use cases. We take the position that the Union's regulatory framework should not foreclose the choice.

\section{Conclusion}

OCPP-AI is offered as an open specification suitable for normative or informative reference in published guidelines and harmonized standards under the EU AI Act. The specification is architecturally compatible with W3C Verifiable Credentials 2.0, architecturally complementary to eIDAS 2.0, GDPR-compatible by structure, and substrate-agnostic with named reference implementations including Bitcoin \texttt{OP\_RETURN} and the European Blockchain Services Infrastructure.

The maintainer (LedgerProof Foundation) is in formation under counsel and does not assert proprietary rights over the specification or commercial encumbrance on its implementation. Stakeholder consultation is invited from the European Commission AI Office, the European AI Board, Member State competent authorities, CEN-CENELEC JTC 21, the IETF SCITT Working Group, and W3C.

\begin{thebibliography}{99}
\bibitem{aiact} European Parliament and Council. \emph{Regulation (EU) 2024/1689 (Artificial Intelligence Act)}. Official Journal of the European Union, 2024. Article 50.
\bibitem{gdpr} European Parliament and Council. \emph{Regulation (EU) 2016/679 (General Data Protection Regulation)}. Official Journal of the European Union, 2016. Articles 4, 17.
\bibitem{eidas2} European Parliament and Council. \emph{Regulation (EU) 2024/1183 amending Regulation (EU) 910/2014 as regards establishing the European Digital Identity Framework (eIDAS 2.0)}. Official Journal of the European Union, 2024.
\bibitem{vc20} World Wide Web Consortium. \emph{Verifiable Credentials Data Model v2.0}. W3C Recommendation, 15 May 2025. \url{https://www.w3.org/TR/vc-data-model-2.0/}
\bibitem{rfc2119} Bradner, S. \emph{Key words for use in RFCs to Indicate Requirement Levels}. RFC 2119, IETF, March 1997.
\bibitem{rfc6234} Eastlake, D., Hansen, T. \emph{US Secure Hash Algorithms (SHA and SHA-based HMAC and HKDF)}. RFC 6234, IETF, May 2011.
\bibitem{rfc6962} Laurie, B., Langley, A., Kasper, E. \emph{Certificate Transparency}. RFC 6962, IETF, June 2013.
\bibitem{rfc8032} Josefsson, S., Liusvaara, I. \emph{Edwards-Curve Digital Signature Algorithm (EdDSA)}. RFC 8032, IETF, January 2017.
\bibitem{rfc8174} Leiba, B. \emph{Ambiguity of Uppercase vs Lowercase in RFC 2119 Key Words}. RFC 8174, IETF, May 2017.
\bibitem{rfc8949} Bormann, C., Hoffman, P. \emph{Concise Binary Object Representation (CBOR)}. RFC 8949, IETF, December 2020.
\bibitem{rfc9162} Laurie, B., Messeri, E., Stradling, R. \emph{Certificate Transparency Version 2.0}. RFC 9162, IETF, December 2021.
\bibitem{scitt-arch} Birkholz, H., et al. \emph{An Architecture for Supply Chain Integrity, Transparency and Trust (SCITT)}. Internet-Draft draft-ietf-scitt-architecture, IETF, work in progress.
\bibitem{dawkins-scitt} Dawkins, V.\ S. \emph{A SCITT Profile for EU AI Act Article 50 Transparency Receipts}. Internet-Draft draft-dawkins-scitt-ai-article50-00, IETF, 25 May 2026.
\bibitem{ebsi-arf} European Commission DG-DIGIT. \emph{European Blockchain Services Infrastructure Architecture Reference Framework v3}. Brussels, 2024.
\bibitem{eif} European Commission. \emph{European Interoperability Framework (EIF) v2.0}. COM(2017) 134 final, Brussels, 2017.
\end{thebibliography}

\end{document}
